\section{Sharpness for two exceptional families}\label{sec:sharpness}

Theorems~A and~B exclude two classes by proving lower bounds on frozen point areas. The
source-cover paper excludes every common power envelope $|G_i|\le HL/i^p$ with $p>q$
\cite[Corollary~4.4]{SourceCovers2026}; for two exceptional families the inverse-square
envelope $p=2$ is the boundary case. This section shows that the frozen gap-area condition
\eqref{eq:GA} cannot exclude that boundary case at moduli comparable to $R$. We construct
admissible geometries with two exceptional families and hop counts $|G_i|\le512L/i^2$
whose frozen point areas have mean at most $10^7/R^2$ over legal phases and columns, at
every modulus, and whose hop tails are at least $1/(32R)$. Hence \eqref{eq:GA} holds on
them whenever $3.2\cdot10^8(M-R)\le\beta R$, for instance at $M=2R+2$ for every
$\beta\ge10^9$ (\cref{thm:sharpness,cor:sharpness-limit}). The geometries carry a single
input. Frozen point areas and hop profiles do not involve the correctness sets, so this is
a statement about what the certificate can see; it constructs no transposer.

\paragraph{The idea.}
Two banks of zigzags, A and~B, carry the root of every column down through bank~A,
across the newer exceptional family $g$ into bank~B, and down through bank~B; the input
head does not move during them. Bank~A turns the column and the depth $p$ at which the
chain leaves it into an offset $o$, one of about $2^\ell$ values when the protected
zigzags have levels up to $\ell$. Bank~B carries every such offset to each of its
zigzags $\ZB_k$. A third bank, C, is swept once by the input head. The older exceptional
family $e$ sends the descent and the first rise of $\ZB_k$ into the second rise of the
bank-C zigzag $\ZC_k$, at relative levels $\vdc(k)\mp o/4$ around a centre given by the
base-$2$ van der Corput sequence. The landing levels of one column, about $4^\ell$ in
number, fill the unit interval up to $O(4^{-\ell})$ (\cref{lem:rows}); a single bank
would give only about $2^\ell$ of them. The chains of bank~C then copy every landing
level into every earlier zigzag of bank~C, so the frozen raw values fill the input range
except for one arc of length $O(L/2^\ell)$ (\cref{lem:dense-gaps}). Landings that fall
outside a second rise are parked elsewhere and only add points.

\begin{definition}[The dyadic sharpness geometries]\label{def:sharp-geometry}
\formal{GapAreaMemo/Z3DGeometry.lean}{SourceCovers.Z3D.Index.geometry,SourceCovers.Z3D.Index.runtime_le}
Fix an integer $\nu\ge3$ and put $d=2^{2^\nu}$, so that $b=b(d)=2^\nu$ and
$L=d^2b=2^{2^{\nu+1}+\nu}$. Each bank has zigzags with indices
$1\le k\le k_\nu=2^{\nu+1}-1$; the index $k$ has level $\ell(k)=\lfloor\log_2k\rfloor\le\nu$
and weight $w_k=4^{-\ell(k)-1}$. Banks A and~B have heights $w_kL$ and bank~C has
heights $8w_kL$; consecutive heights differ by a factor $1$ or~$4$. A zigzag of height
$h$ rises $h$ cells, descends $h$ cells and rises $h$ cells again; every leg moves one
cell per step, except that one rising leg may be slowed by a constant factor. Write
$\vdc$ for the base-$2$ van der Corput sequence, which reverses binary digits:
$\vdc\bigl(\sum_ik_i2^i\bigr)=\sum_ik_i2^{-i-1}$ for digits $k_i\in\{0,1\}$.
\begin{enumerate}[label=(\roman*)]
\item \emph{Bank C} consists of $\ZC_1,\dots,\ZC_{k_\nu}$ in this time order. $\ZC_1$
starts at time~$1$; for $k\ge2$, the first rise of $\ZC_k$ runs during the second rise
of $\ZC_{k-1}$, slowed by the factor $w_{k-1}/w_k$. Let $U_k$ be the trough of $\ZC_k$,
and let $U^*=U_{k_\nu}+8w_{k_\nu}L$ be the end of bank~C.
\item \emph{Bank B} consists of $\ZB_{k_\nu},\dots,\ZB_1$ in this time order, after
bank~C. The second rise of $\ZB_k$ runs during the first rise of $\ZB_{k-1}$, slowed by
the factor $w_{k-1}/w_k$; that of $\ZB_1$ is slowed by~$4$ and takes $L$ steps. Bank~B
ends at the trough $O_B$ of~$\ZB_1$.
\item \emph{Bank A} consists of zigzags $\ZA_k$ laid out like those of bank~B, after the
second rise of $\ZB_1$; it ends at the trough $O_A$ of $\ZA_1$. The outputs are appended
during the second rise of $\ZA_1$, at $t_z=O_A+1+z$ for $z<L$.
\item \emph{Heads.} One work head performs the zigzags of odd index of the three banks in
time order, another those of even index. Each zigzag starts at a fresh cell above all
cells its head has visited before. After
the halting time each of these heads performs padding zigzags; complete families are
formed from hops before the halting time only, so the padding adds no family. Two
further work heads perform one excursion each. The newer family $g$ rises at unit speed
from $O_B+1$ and descends at unit speed during bank~A, reaching level $i$ at time
$O_A-i$, for $i<h_g=2\sum_kw_kL$. The older family $e$ descends at unit speed during
bank~B, reaching level $l$ at time $O_B-l$, and rises earlier. Level $0$ departs at
time $0$; for $l\ge1$, the level that reaches the bank-B time at offset $0\le r<2w_kL$
from the start of $\ZB_k$ departs at the landing time
\[
 U_k+\psi_k\bigl(\lfloor4w_kL\,\vdc(k)\rfloor+w_kL-r\bigr),
\]
where $\psi_k(z)=2z$ for $0\le z\le4w_kL$ (a landing in the second rise of $\ZC_k$),
$\psi_k(z)=z$ for $z<0$ and $\psi_k(z)=4w_kL+z$ for $z>4w_kL$ (parking in the descents
of $\ZC_k$ and $\ZC_{k+1}$, or, for $k=k_\nu$, shortly after bank~C).
\item \emph{Input head.} It sweeps once across bank~C:
$x(t)=\min(\lfloor(t-1)/\sigma\rfloor,L-1)$ with $\sigma=\lfloor(U^*-1)/L\rfloor$.
\item \emph{Computation.} A single input and a one-point record alphabet.
\end{enumerate}
The original core consists of the upper and lower families of the $3k_\nu$ bank
zigzags; the exceptional families are $e$ and~$g$; and the bound on all other families
is $a_0=0$. We write $\Gamma_\nu$ for this geometry.
\end{definition}

The Lean namespace is \texttt{Z3D}, for three banks of zigzags with dyadic parameters. The
parameters are \lean{SourceCovers.Z3D.Index}, the heads
\lean{SourceCovers.Z3D.Index.heads}, the raw path and the append times
\lean{SourceCovers.Z3D.Index.raw} and \lean{SourceCovers.Z3D.Index.out}, the two
exceptional schedules \lean{SourceCovers.Z3D.Index.eSchedule} and
\lean{SourceCovers.Z3D.Index.gSchedule}, and the core
\lean{SourceCovers.Z3D.Index.core}. Each zigzag splices two scheduled excursions of the
shared library.

\begin{lemma}[Admissibility]\label{lem:sharp-admissible}
\formal{GapAreaMemo/Z3DFamilies.lean}{SourceCovers.Z3D.Index.family_mem,SourceCovers.Z3D.Index.core_disjoint,SourceCovers.Z3D.Index.exceptions_card,SourceCovers.Z3D.Index.outside_none}
For every $\beta\ge0$, $\Gamma_\nu$ is an admissible geometry with $c=128$, $N=3$,
$q=2$ and $a_0=0$. Every complete family is a core family, $e$ or~$g$; the core has a
two-class span partition and the full family system a three-class one.
\end{lemma}

\begin{outline}
The walk of each of the first two heads is a sequence of zigzags at fresh cells, and its
complete families are exactly the upper families of its zigzags (the returns of the
first rise and the descent, owned at the peak) and the lower ones (the returns of the
descent and the second rise, owned at the trough); each exceptional excursion is one
native family. Hence every family is a bank family or an exception. Upper families of
distinct bank zigzags have disjoint spans, and so do lower families; colouring the core
by upper and lower gives its two classes (\lean{SourceCovers.Z3D.Index.core_coloring}).
The spans of $e$ and $g$ are disjoint, and they form the third class
(\lean{SourceCovers.Z3D.Index.full_coloring}). The banks and the output phase take at
most $128L$ steps.
\end{outline}

\begin{lemma}[Inverse-square sizes and protection]\label{lem:sharp-sizes}
\formal{GapAreaMemo/Z3DProtected.lean}{SourceCovers.Z3D.Index.famA_card,SourceCovers.Z3D.Index.famB_card,SourceCovers.Z3D.Index.famC_card,SourceCovers.Z3D.Index.card_big_le,SourceCovers.Z3D.Index.hop_envelope,SourceCovers.Z3D.Index.famA_protected,SourceCovers.Z3D.Index.famB_protected,SourceCovers.Z3D.Index.famC_protected,SourceCovers.SortedCore.prefix_retained}
In $\Gamma_\nu$, the upper and lower families of $\ZA_k$ and of $\ZB_k$ have $w_kL$ hops
each, and those of $\ZC_k$ have $8w_kL$ hops. For every $\ell\ge0$, at most
$16\cdot2^\ell-6$ core families have at least $L\,4^{-\ell-1}$ hops. Hence, in every order
$G_1,G_2,\dots$ of the core by nonincreasing hop count:
\begin{enumerate}[label=(\alph*)]
\item $|G_i|\le512L/i^2$ for every $i$;
\item if $16\cdot2^\ell\le R$, the first $R$ families contain the upper and lower families
of every $\ZA_k$ and $\ZB_k$ with $\ell(k)\le\ell$ and of every $\ZC_k$ with
$\ell(k)\le\ell+1$; and for all $J\ge R$ and $M$, every family among the first $R$
survives at every legal phase.
\end{enumerate}
\end{lemma}

\begin{outline}
Each core family is the upper or lower family of one bank zigzag, and its hops are the
returns of one pair of legs, one per cell. The families with at least $L4^{-\ell-1}$ hops
belong to zigzags $\ZA_k$, $\ZB_k$ with $\ell(k)\le\ell$ and $\ZC_k$ with
$\ell(k)\le\ell+1$, at most $4(2^{\ell+1}-1)+2(2^{\ell+2}-1)$ families. (a) If
$L4^{-\ell-1}\le|G_i|<L4^{-\ell}$, then $G_1,\dots,G_i$ all have at least $L4^{-\ell-1}$
hops, so $i<16\cdot2^\ell$; and every family has at most $2L$ hops. (b) If a family with
at least $L4^{-\ell-1}$ hops stood at position $p>R$, the $p$ distinct families up to it
would all have that many hops, more than $16\cdot2^\ell$. A protected family survives at a
legal phase \cite[Lemma~2.6]{SourceCovers2026}.
\end{outline}

\begin{lemma}[Van der Corput rows with dyadic offsets]\label{lem:rows}
\formal{GapAreaMemo/Covering.lean}{SourceCovers.Covering.vdC,SourceCovers.Covering.vdC_hits,SourceCovers.Covering.w,SourceCovers.Covering.t,SourceCovers.Covering.t_sub_succ,SourceCovers.Covering.offsets,SourceCovers.Covering.offsets_dense,SourceCovers.Covering.row_cover_quarter,SourceCovers.Covering.rowPoints_dense_of_le}
For $p\ge1$ let $w_p$ be as in \cref{def:sharp-geometry}, put
$t_p=(2^{\ell(p)+1}-p)2^{-2\ell(p)-1}+2^{-\ell(p)-1}$, so that $t_1=1$ and
$t_p-t_{p+1}=2w_p$, and for $D\ge1$ and rational $\theta\in[0,1]$ let
$O_D(\theta)=\{t_p-\theta w_p,\ t_p-(2-\theta)w_p:1\le p\le D\}$. Let $n\ge6$ and
$2^n\le16D$. Then for every $k_0\ge0$ the points $\vdc(k)\pm o/4$ with $k_0\le k<k_0+2^n$
and $o\in O_D(\theta)$ meet every interval $[x,y]\subseteq[0,1]$ with rational endpoints
and $y-x\ge1152\cdot4^{-n}$.
\end{lemma}

\begin{outline}
Any $2^n$ consecutive values of $\vdc$ meet every dyadic interval of length $2^{-n}$,
hence every interval of length $2^{1-n}$ in $[0,1]$. The offsets lie in $(0,1]$, and above
every $r\in[t_D,1/2]$ some offset lies in $[r,r+8r^2]$, because consecutive values $t_p$
differ by $2w_p$ and $w_p\le r^2$ when $t_{p+1}\le r$. Given $[x,y]$, choose a centre
$\vdc(k)$ at a distance of order $2^{-n}$ from one end of the interval and an offset close
to four times that distance: the row point $\vdc(k)\pm o/4$ then lies in $[x,y]$. These
covering properties are elementary, and we claim no novelty for them.
\end{outline}

\begin{lemma}[Point area from dense values]\label{lem:dense-gaps}
\formal{GapAreaMemo/GapBound.lean}{SourceCovers.GapBound.window_cover,SourceCovers.GapBound.pointArea_le_of_dense'}
\begin{enumerate}[label=(\alph*)]
\item Let $\alpha_0,\dots,\alpha_n$ and $w'_0,\dots,w'_{n-1}$ be natural numbers with
$\alpha_s+w'_s\le\alpha_{s+1}$, and let $\Lambda\ge1$ be at least every $w'_s$ and every
$w'_s+w'_{s+1}$. If every window $[z,z+w'_s]\subseteq[\alpha_s,\alpha_{s+1}]$ meets a set
$S$ of natural numbers, then every window $[z,z+\Lambda]\subseteq[\alpha_0,\alpha_n]$
meets~$S$.
\item Let $S\subseteq\Omega$ be nonempty with $L-1\in S$, and let $\gamma\le X\le L$ be
integers such that every window $[u,u+\gamma]\subseteq[0,X]$ meets $S$. Then
\[
 \gaparea(S)\le\frac{(\gamma+1)L+(L-X+\gamma+1)^2}{2L^2}.
\]
\end{enumerate}
\end{lemma}

\begin{outline}
(a) Let $\alpha_s\le z<\alpha_{s+1}$. If $z+w'_s\le\alpha_{s+1}$, use the window
$[z,z+w'_s]$; otherwise use $[\alpha_{s+1},\alpha_{s+1}+w'_{s+1}]$, which lies in
$[z,z+\Lambda]$. (b) A gap starting at a point $u$ with $u+\gamma+1\le X$ has length at
most $\gamma+1$. The other gaps that do not wrap around lie in $[X-\gamma,L-1]$ and have
total length at most $L-X+\gamma+1$, and the wraparound gap ends at a point at
most~$\gamma$. Bound the sum of the squared gaps by the largest gap of each kind times the
total length of that kind.
\end{outline}

\begin{theorem}[Theorem Z: frozen gap areas of the sharpness geometries]\label{thm:sharpness}
\formal{GapAreaMemo/Z3DSharp.lean}{SourceCovers.Z3D.Index.phase_le,SourceCovers.Z3D.Index.mean_le,SourceCovers.Z3D.Index.gap_area_le,SourceCovers.Z3D.Index.hop_tail,SourceCovers.Z3D.Index.gap_area_condition,SourceCovers.Z3D.Index.gap_area_condition_default}
Let $\nu\ge3$ and $\beta\ge0$, and fix an order $G_1,G_2,\dots$ of the core of
$\Gamma_\nu$ by nonincreasing hop count.
\begin{enumerate}[label=(\alph*)]
\item If $1\le R\le2^\nu$, $M\ge R+1$ and $J\ge R$, then
$\meanarea_{J,R,M}\le10^7/R^2$.
\item If $R\ge1$, $4R\le2^\nu$ and $J\ge64R$, then $\sum_{R<i\le J}|G_i|\ge L/(32R)$.
\item Hence, if $R\ge1$, $4R\le2^\nu$, $J\ge64R$, $M\ge R+1$ and
$3.2\cdot10^8(M-R)\le\beta R$, then
\[
 \meanarea_{J,R,M}\le\frac{\beta}{M-R}\sum_{R<i\le J}\frac{|G_i|}L.
\]
In particular this holds for $M=2R+2$ whenever $\beta\ge10^9$.
\end{enumerate}
\end{theorem}

\begin{outline}
(a) For $R^2\le2\cdot10^7$ the bound is at least $1/2$. Otherwise choose $\ell\ge8$ with
$16\cdot2^\ell\le R<32\cdot2^\ell$. By \cref{lem:sharp-sizes}(b), at every legal phase the
orbit families contain the families of all bank zigzags of indices below $2^{\ell+1}$
and both exceptions. Fix a column $j\ge1$ and put $\theta=j/d$.
\begin{enumerate}[label=(\roman*)]
\item \emph{Bank A.} The root $s_j$ lies in the second rise of $\ZA_1$ at level
$\theta w_1L$. The lower family of $\ZA_1$ maps it into the descent, the upper family into
the first rise, which runs during the second rise of $\ZA_2$, and so on: $\ZA_p$ is
reached at the exact level $\theta w_pL$, in its descent and in its first rise
(\lean{SourceCovers.Z3D.Index.late_chain}).
\item \emph{Across $g$.} These times lie in the descent of $g$, which maps them into the
second rise of $\ZB_1$ at relative level $1-o$, with $o=t_p-\theta w_p$ from the
descent and $o=t_p-(2-\theta)w_p$ from the first rise, because
$2\sum_{i<p}w_i=1-t_p$.
\item \emph{Bank B and $e$.} The chain of bank~B reaches $\ZB_k$ at a level within $3$ of
$(1-o)w_kL$, the levels being rounded down. The family $e$ maps the descent time of
$\ZB_k$ into the second rise of $\ZC_k$ at a height $2\mu$ whose relative level
$(2\mu+1)/(8w_kL)$ lies within $1/(w_kL)$ of $\vdc(k)-o/4$, and the first-rise time to
one within $1/(w_kL)$ of $\vdc(k)+o/4$ (\lean{SourceCovers.Z3D.Index.landing}).
\item \emph{Density.} By \cref{lem:rows} with $k_0=2^\ell$, $n=\ell$ and $D=2^\ell$,
these relative levels, for $2^\ell\le k<2^{\ell+1}$, meet every window inside
$[2/(w_kL),1-2/(w_kL)]$ of relative length about $1152\cdot4^{-\ell}$
(\lean{SourceCovers.Z3D.Index.levels_window}).
\item \emph{Bank C.} The chain of bank~C carries a landing at height $\eta$ in $\ZC_k$ to
every $\ZC_i$ with $i\le k$, at the height $\eta'$ with
$(\eta'+1)/(8w_iL)=(\eta+1)/(8w_kL)$, in its second rise, its descent and its first rise
(\lean{SourceCovers.Z3D.Index.early_chain}). Hence every window of relative length about
$1152\cdot4^{-\ell}$ inside the first rise or the descent of each of
$\ZC_1,\dots,\ZC_{2^\ell}$ meets the orbit, and by \cref{lem:dense-gaps}(a) so does every
time window of length $4608L4^{-\ell}+O(4^\ell)$ that ends by the end of the first rise
of $\ZC_{2^\ell}$ (\lean{SourceCovers.Z3D.Index.time_cover}).
\item \emph{Raw values.} Since $x$ advances one cell every $\sigma\ge1$ steps, the frozen
raw set meets every window of that length inside $[0,X]$. Here
$L-X\le40L/2^\ell+O(L4^{-\ell})$, because the rest of bank~C lasts at most
$40\sum_{k>2^\ell}w_kL+16w_{2^\ell}L$ steps and $\sigma\ge2$; and $L-1=x(s_j)$ lies in the
set. \Cref{lem:dense-gaps}(b) gives $\gaparea(S_\rho(s_j))\le4010\cdot4^{-\ell}$
(\lean{SourceCovers.Z3D.Index.phase_bound}).
\end{enumerate}
None of this depends on $M$. Averaging over legal phases and columns, with the column
$j=0$ bounded by $1/2$, gives at most $4011\cdot4^{-\ell}<4.2\cdot10^6/R^2$. (b) Let
$2^r$ be the power of two with $2R<2^r\le4R$. The upper families of the zigzags $\ZA_k$
with $2^r\le k<2^{r+1}$ are $2^r$ distinct families of $L4^{-r-1}$ hops each. Fewer than
$16\cdot2^r\le J$ core families have that many hops, so all of them lie among the first $J$
(\cref{lem:sharp-sizes}); at most $R$ of them lie among the first $R$; hence
$\sum_{R<i\le J}|G_i|\ge(2^r-R)L4^{-r-1}\ge L/(32R)$. (c) Combine (a) and~(b):
$10^7/R^2\le\beta/(32R(M-R))$; for $M=2R+2$ use $M-R\le3R$.
\end{outline}

\begin{corollary}[Limit form]\label{cor:sharpness-limit}
\formal{GapAreaMemo/Z3DLimit.lean}{SourceCovers.Z3D.seq,SourceCovers.Z3D.limit_form,SourceCovers.Z3D.limit_GA,SourceCovers.Z3D.limit_GA_zero,SourceCovers.Z3D.limit_GA_default}
Consider the sequence $(\Gamma_{\nu+3})_{\nu\ge0}$ with any orders of its cores by
nonincreasing hop count, and any strictly increasing subsequence along which the
normalized hop profiles $|G_i|/L$ converge to limits $h_i$. Define $\limarea_{R,M}$ and
$\tau_R=\sum_{i>R}h_i$ as in \cref{def:gap-area-means} and \cref{sec:preliminaries}.
Then $\limarea_{R,M}\le10^7/R^2$ and $\tau_R\ge1/(32R)$ for all $R\ge1$ and $M\ge R+1$.
Hence
\[
 \limarea_{R,M}\le\frac{\beta\tau_R}{M-R}
\]
whenever $R\ge1$, $M\ge R+1$ and $3.2\cdot10^8(M-R)\le\beta R$, and whenever $R=0$ and
$2\le M\le\beta/16$. In particular \eqref{eq:GA} holds at $M=2R+2$ for every $R\ge0$
when $\beta\ge10^9$.
\end{corollary}

\begin{outline}
For fixed $J\ge R$, the geometries of the subsequence eventually satisfy $R\le2^\nu$, and
\cref{thm:sharpness}(a) bounds the inner upper limit; the outer one preserves the bound.
The finite profile masses are at most $128$, so the limit profile is nonnegative and
summable; \cref{thm:sharpness}(b) with $J=64R$ passes to the limit, and
$\tau_R\ge\sum_{R<i\le64R}h_i$. For $R=0$, $\limarea_{0,M}\le1/2$ and
$\tau_0\ge\tau_1\ge1/32$.
\end{outline}

\paragraph{Why it matters.}
For an impossibility proof, Theorem~Z closes one route: the frozen gap-area condition at
moduli up to $R+\beta R/(3.2\cdot10^8)$ cannot exclude the two-exception inverse-square
class. An exclusion of that class must use larger moduli, the remaining necessary
conditions of~\cite{SourceCovers2026}, or input specific to transposition
(\cref{prob:transposition-input}). For a construction, the geometries show how two
exceptional families can meet \eqref{eq:GA}: the older family lands the chains of the
protected banks on both sides of well-spread centres in a bank that the input head
sweeps, and the chains of that bank merge the landings of different zigzags.

The sequence of \cref{cor:sharpness-limit} belongs to the class of
\cref{prob:two-exception} for all $c\ge128$, $N\ge3$, $H\ge512$ and $\beta>0$, with
$a_0=0$, and \eqref{eq:GA} holds along every joint profile at every modulus allowed in
\cref{cor:sharpness-limit}. The argument of \cref{rem:reciprocal-comparison} does not
use near-fullness; applied to this sequence (proof sketch only), it shows that the
weighted reciprocal-cover bound \cite[(4.8)]{SourceCovers2026} holds at the same moduli.
By \cref{lem:funnel-area}(c), \cref{thm:sharpness} also implies that for large $\nu$ the
geometries have no funnel of a fixed depth over a raw interval of fixed positive
relative length; and they are not single-sweep terminal-landing geometries, since the
input head moves on the lower regions of the bank-C families. These two observations
are elementary and are not formalized. The constants are not optimized. Nothing here
concerns larger moduli, the necessary complete-observable movement and global
continuation density of \cite[Corollaries~4.8 and~5.3]{SourceCovers2026}, or the
existence of a transposer with such geometries.
